\scanpage{14}

\begin{center}
U N I V E R S I T Y\\
OF SCIENCE AND TECHNOLOGY\\
\uline{OF LANGUEDOC}

Student Records Office\\
MA/GJ\\
Distribution : List A\\
Number: 4680
\end{center}

Montpellier, 20 October 1983

\begin{center}
\uline{NOTICE OF DEFENCE OF A THIRD-CYCLE DOCTORAL THESIS}

\uline{EVOLUTION AND BIOSYSTEMATICS}
\end{center}

\uline{Surname and given name of the candidate} : CHAKROUN Saïda

\uline{Title of the thesis} : THE SEASONAL VARIATIONS IN THE SECONDARY FORMATIONS (WOOD, PHLOEM, CAMBIUM) OF THE TRUNK OF THREE MEDITERRANEAN SPECIES, THE HOLM OAK, THE OLIVE TREE, AND THE HACKBERRY : A HISTOPHYSIOLOGICAL STUDY.

\uline{Scheduled date, time and place of the defence} : 28 October 1983 at 9:30\\
Lecture Room number 9 - University of Science and Technology of Languedoc

\uline{Name of the thesis supervisors} : Misters Charles HEBANT and BONNET

\uline{Designation of the laboratory where the thesis is prepared} : Laboratory of Plant Biology

\uline{Summary of the thesis} :

This study concerns the anatomy, histology and histochemistry of the secondary conducting tissues of \textit{Olea europaea}, of \textit{Quercus ilex}, and of \textit{Celtis australis}.

The wood and phloem of the Oak and the Olive tree have been the subject of an in-depth structural study.

The functioning of the cambial zone (as well as the formation of its derivatives, wood and phloem) and the evolution of the starch reserves of the trunk of these three species undergo considerable variations with the seasonal rhythm, and in the case of the Olive tree, with the alternations in production.

In the Oak and the Hackberry, all the wood and phloem is differentiated in spring and summer. In the Olive tree, cambial reactivation in spring is more or less early depending on the fruit load ; furthermore, only trees with low production show a resumption of cambial activity in autumn, the latter produces only wood.

In the three species, the formation of the phloem is initiated later than in the case of the wood. In the two evergreen species, the zone of conducting phloem never disappears, on the other hand, in the Hackberry, at the end of autumn and during winter, all the sieve tubes become non-conducting by accumulating callose.

In parallel with cambial reactivation, a hydrolysis of starch is noted. Amylogenesis is more or less early depending on the species observed, and in the case of the Olive tree, depending on the fruit load. A competition between vegetative growth and the reproductive development of the Olive tree is observed.

Throughout the year, the companion cells and the vascular rays of the Oak and the Olive tree possess high phosphohydrolase and respiratory activity.

\uline{Approximate duration of the research leading to the thesis} : 3 years

\uline{Proposed composition of the examining committee} :

\begin{tabular}{ll}
Misters & D'AUZAC J.\\
& BONNET A.L.\\
& KAUMAN W.\\
& NABLI M.A\\
& THALER L.\\
& VILLEMUR P.
\end{tabular}

\begin{center}
For the Secretary General,\\
the Principal Officer for School\\
and University Administration,

\unsure{A. Gentil}

A. GENTIL
\end{center}

\margin{13}
